\usepackage{graphicx}  % Necesar pentru \rotatebox (rotirea textului)
\usepackage{multirow}  % Necesar pentru \multirow (celule unite pe verticală)
\usepackage{array}
\section{Experiment 8: 2D Slice-Based Approach (Transfer Learning)}

\subsection{Rationale and Hypothesis}
While 3D CNNs capture volumetric spatial context, they are computationally intensive and suffer from a lack of large-scale pre-trained models comparable to those in general computer vision. 

Experiment 8 tests the hypothesis that a \textit{2D Slice-Based approach} can outperform 3D models by leveraging:
\begin{itemize}
    \item \textbf{Stronger Transfer Learning:} Utilizing architectures (EfficientNetB0) pre-trained on millions of 2D images (ImageNet).
    \item \textbf{Data Efficiency:} Reducing the 3D problem to 2D simplifies the feature space, potentially reducing the risk of overfitting on our small cohort ($N=87$).
\end{itemize}

\subsection{Methodology}

\subsubsection{Slice Extraction Strategy}
To convert the 3D CT volumes into 2D inputs, we employed a "Max-Tumor-Area" slicing strategy:
\begin{enumerate}
    \item \textbf{ROI Identification:} For each patient, the lung mask was used to identify the bounding box of the tumor.
    \item \textbf{Slice Selection:} We iterated through the axial slices of the cropped volume and selected the single slice with the largest tumor surface area (maximum number of mask pixels).
    \item \textbf{Pseudo-Coloring:} Since standard 2D models expect 3 input channels (RGB), the selected grayscale CT slice was replicated 3 times to fit the input tensor shape $(224, 224, 3)$.
\end{enumerate}

\subsubsection{Model Architecture}
We utilized a standard 2D Convolutional Neural Network architecture compatible with ImageNet weights.
\begin{itemize}
    \item \textbf{Backbone:} EfficientNetB0 (pre-trained on ImageNet).
    \item \textbf{Adaptation:} The top (classification) layers were removed and replaced with:
    \begin{itemize}
        \item GlobalAveragePooling2D.
        \item Dense Layer (128 units, ReLU).
        \item Dropout (0.5).
        \item Dual Output Heads: \texttt{Dense(1, Sigmoid)} for TTF1 and \texttt{Dense(1, Sigmoid)} for CK7.
    \end{itemize}
\end{itemize}

\subsection{Experimental Results}

\subsubsection{Test Set Performance}
The 2D model was evaluated on the same stratified test set.

\begin{figure}[h]
    \centering
    % --- Matrice TTF1 ---
    \begin{minipage}{0.45\textwidth}
        \centering
        \textbf{a) TTF1 Prediction} \\ \vspace{0.3cm}
        \renewcommand{\arraystretch}{1.5} % Mărește spațiul vertical
        \begin{tabular}{cc|c|c|}
            & \multicolumn{1}{c}{} & \multicolumn{2}{c}{\textbf{Predicted}} \\
            & \multicolumn{1}{c}{} & \multicolumn{1}{c}{Negative (0)} & \multicolumn{1}{c}{Positive (1)} \\ \cline{3-4}
            \multirow{2}{*}{\rotatebox[origin=c]{90}{\textbf{Actual}}} 
            & Neg (0) & \textbf{2} & 1 \\ \cline{3-4}
            & Pos (1) & 1 & \textbf{5} \\ \cline{3-4}
        \end{tabular}
    \end{minipage}
    \hfill
    % --- Matrice CK7 ---
    \begin{minipage}{0.45\textwidth}
        \centering
        \textbf{b) CK7 Prediction} \\ \vspace{0.3cm}
        \renewcommand{\arraystretch}{1.5}
        \begin{tabular}{cc|c|c|}
            & \multicolumn{1}{c}{} & \multicolumn{2}{c}{\textbf{Predicted}} \\
            & \multicolumn{1}{c}{} & \multicolumn{1}{c}{Negative (0)} & \multicolumn{1}{c}{Positive (1)} \\ \cline{3-4}
            \multirow{2}{*}{\rotatebox[origin=c]{90}{\textbf{Actual}}} 
            & Neg (0) & \textbf{5} & 0 \\ \cline{3-4}
            & Pos (1) & 2 & \textbf{2} \\ \cline{3-4}
        \end{tabular}
    \end{minipage}
    
    \vspace{0.5cm}
    \caption{Confusion Matrices for the 2D Slice-Based Model on the Test Set.}
    \label{fig:conf_matrix_exp8}
\end{figure}

\subsubsection{Comparison with 3D Approaches}
\begin{itemize}
    \item \textbf{Pros:} The 2D model showed [High/Low] sensitivity to texture, likely due to the robust filters learned from ImageNet.
    \item \textbf{Cons:} By discarding the volumetric information (the "z" dimension), the model may have lost 3D morphological features (e.g., sphericity, spiculated margins in 3D) that are relevant for diagnosis.
\end{itemize}

\subsection{Conclusion of Experiment 8}
This experiment demonstrates the trade-off between dimensionality and data availability. While 2D transfer learning offers a powerful starting point, the reduction of a 3D tumor to a single 2D cross-section simplifies the biological reality. The results suggest that [Conclusion based on results: e.g., "a hybrid approach or multi-slice 2.5D approach might be needed to capture both texture and volume"].